\section{Certified inexact oracles}
\label{sec:inexact}

The regridding argument survives certified errors in local optimization,
separator slopes, and incumbent evaluation.  Local errors accumulate along
the selected configuration, once per bag.  They do not accumulate over all
the convex problems solved to construct the certificate.  Throughout this
section, the models satisfy \eqref{eq:model-contract}, the bag gradients
are $M$-Lipschitz, and quadratic growth holds with constant $g>0$ as in
\eqref{eq:model-growth}.  We use $C_0,B_0,\eta,C$ from
\eqref{eq:regridding-constants}, with the model error constant denoted by
$A$.

\subsection{Certified local problems and selected residuals}

Fix one stage and one exactly represented slope $\widehat\lambda_t$ per
nonroot separator.  The partitions are fixed during the bottom-up pass.
Once the child intercepts $\widehat\beta_{u,D'}$ are available, put
\[
 b_{u,B}=\min\{\widehat\beta_{u,D'}:
                 D'\in\mathcal P_u,\ D'\cap B_{S_u}\ne\varnothing\}
\]
for each parent leaf $B$, and store a minimizing child cell.  For a
nonroot bag, each touching pair $(B,D)$ has the nonempty compact domain
$C_{t,B,D}=\{v\in B:v_{S_t}\in D\}$.  Its local objective is
\begin{equation}
 q_{t,B,D}(v)=m_{t,B}(v)
       +\sum_{u\in\operatorname{ch}(t)}
                  [\widehat\lambda_u^T v_{S_u}+b_{u,B}]
       -\widehat\lambda_t^T v_{S_t}.
 \label{eq:inexact-local-objective}
\end{equation}
At the root the domain is $B$ and the final slope term is absent.  A
certified local oracle returns lower and upper numbers $\ell_{t,B,D}$,
$u_{t,B,D}$ and a point $v_{t,B,D}$ satisfying
\begin{equation}
 v_{t,B,D}\in C_{t,B,D},\qquad
 \ell_{t,B,D}\le\min_{C_{t,B,D}}q_{t,B,D},\qquad
 q_{t,B,D}(v_{t,B,D})\le u_{t,B,D}
                 \le\ell_{t,B,D}+\delta_{t,j}.
 \label{eq:inexact-oracle}
\end{equation}
Here $\delta_{t,j}\ge0$ is a budget shared by the tasks of bag $t$ at
stage $j$.  Feasibility and the inequalities are exact statements, not
numerical tolerances.  Rational bounds and points are allowed; other
exact representations are also allowed when the required comparisons
and arithmetic are available.

Set
\begin{equation}
 \widehat\beta_{t,D}
     =\min_{B:B_{S_t}\cap D\ne\varnothing}\ell_{t,B,D},
 \qquad
 \LB_j=\min_{B\in\mathcal L_r}\ell_{r,B},
 \label{eq:inexact-dp}
\end{equation}
and store a minimizing leaf and its oracle point for each table entry.
Backtracking starts at a minimizing root leaf.  In each selected leaf it
uses the stored point and the child cells minimizing $b_{u,B}$, then
continues with those child entries.  All intersections, including faces,
are included in this construction.

\begin{proposition}[Validity and exact selected residuals]
 \label{prop:inexact-dp}
The preceding construction gives a certificate in
\cref{def:model-certificate}, with cell functions
$\widehat l_{t,D}(s)=\widehat\lambda_t^Ts+\widehat\beta_{t,D}$ and
child functions $\widehat\psi_{u,B}(s)=\widehat\lambda_u^Ts+b_{u,B}$.
In particular, $\LB_j\le f^*$ for arbitrary slopes and nonnegative
budgets.  Backtracking gives a configuration, whose value is
\begin{equation}
 \widehat\Phi=\sum_t m_{t,B_t}(z^t)
       +\sum_{t\ne r}\widehat\lambda_t^T
                        (z^{\pi(t)}_{S_t}-z^t_{S_t}).
 \label{eq:inexact-configuration}
\end{equation}
If $\ell_t$ and $q_t$ denote its selected local bound and objective, then
\begin{equation}
 \widehat\Phi-\LB_j=\sum_t\rho_t,
 \qquad \rho_t=q_t(z^t)-\ell_t\in[0,\delta_{t,j}].
 \label{eq:inexact-residual}
\end{equation}
Consequently,
\begin{equation}
 0\le\widehat\Phi-\LB_j\le\sum_t\delta_{t,j}.
 \label{eq:inexact-backtrack}
\end{equation}
\end{proposition}

\begin{proof}
For every child cell meeting $B_{S_u}$, the definition gives
$b_{u,B}\le\widehat\beta_{u,D'}$.  Thus the child function lies below
that cell function throughout their intersection, as required by the
child-minorant inequality.  For an own-separator pair, every feasible
$v$ satisfies
$q_{t,B,D}(v)\ge\ell_{t,B,D}\ge\widehat\beta_{t,D}$.
Adding the incoming affine term gives the local certificate inequality.
The root inequality follows from $q_{r,B}(v)\ge\ell_{r,B}\ge\LB_j$.
Certificate validity now follows from \cref{prop:model-validity}; no
maximality of the intercepts is needed.

Each backtracked point belongs to its selected leaf and own cell, and
each selected child cell meets the selected parent projection.  These
are precisely the configuration conditions in \eqref{eq:model-config}.
For every selected nonroot task,
$\ell_t=\widehat\beta_{t,D_t}$, while the selected root bound equals
$\LB_j$.  Also $b_{t,B_{\pi(t)}}
=\widehat\beta_{t,D_t}$ by the stored child choice.  Summing
\eqref{eq:inexact-local-objective} over selected bags therefore gives
\[
 \sum_t q_t(z^t)=\widehat\Phi+\sum_{t\ne r}\widehat\beta_{t,D_t},
 \qquad
 \sum_t\ell_t=\LB_j+\sum_{t\ne r}\widehat\beta_{t,D_t}.
\]
Subtracting proves the identity.  Finally,
$\ell_t\le q_t(z^t)\le u_t\le\ell_t+\delta_{t,j}$ by
\eqref{eq:inexact-oracle}.  This proves the residual intervals and their
sum bound.
\end{proof}

Thus errors already inherited through child intercepts are counted once
by \eqref{eq:inexact-residual}.  No factor for the number of leaves,
cells, touching pairs, or oracle calls enters this bound.  In an
implementation, the inherited constants $b_{u,B}$ can be omitted from
the numerical optimization objective and added back exactly to both
returned bounds.

\subsection{Approximate slopes and bag gradients}

For a configuration at center $c$, use the regridding quantities
\[
 b_t=\width(B_t),\quad d_t=\width(D_t),\quad
 Q=\sum_t b_t^2+\sum_{t\ne r}d_t^2,\quad R=\norm{x-c},
\]
where $x_i=z_i^{\topbag(i)}$.  Let $\Phi$ be its value with the exact
current-center slopes $\lambda_t(c)$ of \eqref{eq:model-slope}, and set
\[
 C_1=2p\max\{k-1,1\},\qquad
 \nu^2=\sum_{t\ne r}\norm{\widehat\lambda_t-\lambda_t(c)}^2.
\]

\begin{lemma}[Slope perturbation]
 \label{lem:inexact-slopes}
Every configuration satisfies
\begin{equation}
 H:=\sum_{t\ne r}\norm{z^{\pi(t)}_{S_t}-z^t_{S_t}}^2
       \le C_1Q,\qquad
 \abs{\widehat\Phi-\Phi}\le\nu\sqrt{C_1Q}.
 \label{eq:inexact-edge-incidence}
\end{equation}
Suppose the partitions are graded at width $h_j$, the first condition
of \eqref{eq:regridding-theta} holds, and
$\nu\le\zeta\sqrt N h_j$.  Then
\begin{equation}
 \abs{\widehat\Phi-\Phi}
       \le\frac\eta4R^2+K_sNh_j^2,
 \qquad
 K_s=\sqrt{12C_1}\,\zeta
              +\frac{12kC_1\zeta^2\theta^2}{\eta}.
 \label{eq:inexact-slope-perturbation}
\end{equation}
\end{lemma}

\begin{proof}
For a separator coordinate, the child point lies in $D_t$ and the parent
point lies in $B_{\pi(t)}$.  Since their projections
meet, their coordinate difference is at most
$b_{\pi(t)}+d_t$.  Hence
\[
 H\le2\sum_t\sum_{u\in\operatorname{ch}(t)}|S_u|b_t^2
            +2\sum_{u\ne r}|S_u|d_u^2.
\]
Each coordinate of $V_t$ occurs in at most $k-1$ child bags, since the
total number of its occurrence bags is at most $k$.  Therefore
$\sum_{u\in\operatorname{ch}(t)}|S_u|\le(k-1)|V_t|$ and
\[
 H\le2p(k-1)\sum_t b_t^2+2p\sum_{t\ne r}d_t^2\le C_1Q.
\]
Cauchy--Schwarz applied to the slope-error terms proves the second
inequality in \eqref{eq:inexact-edge-incidence}.  By
\cref{lem:regridding-grading},
$Q\le12Nh_j^2+12k\theta^2R^2$, so
\[
 \abs{\widehat\Phi-\Phi}
 \le\sqrt{12C_1}\,\zeta Nh_j^2
       +\sqrt{12kC_1}\,\zeta\theta\sqrt N h_jR.
\]
Young's inequality bounds the second term by
$\eta R^2/4+12kC_1\zeta^2\theta^2Nh_j^2/\eta$, giving the claim.
When $k=1$, every separator is empty, so $H=\nu=0$ directly.
\end{proof}

\begin{lemma}[From bag-gradient errors to slope errors]
 \label{lem:inexact-gradients}
Let $e_t=\widehat g_t-\nabla a_t(c_{V_t})$ and form
$\widehat\lambda_t$ by the same subtree sums as in
\eqref{eq:model-slope}, with $\widehat g_t$ replacing the exact bag
gradient.  If the sums are exact, then
\begin{equation}
 \nu^2\le G\sum_t\norm{e_t}^2,
 \qquad G=\frac{k(k-1)}2.
 \label{eq:inexact-gradient-incidence}
\end{equation}
In particular, $\norm{e_t}\le\gamma h_j$ for each bag ensures the
budget of \cref{lem:inexact-slopes} with $\zeta=\sqrt G\,\gamma$.
Componentwise error at most $\gamma h_j/\sqrt p$ is sufficient.
\end{lemma}

\begin{proof}
Fix a coordinate $i$.  Its occurrence bags form a rooted tree $T_i$ of
size $m_i\le k$.  For every edge below its top bag, the corresponding
slope error is the sum of the errors $e_{t,i}$ in the occurrence bags
below that edge.  Write this linear map as $A_i$.  Its entries are zero
or one, and a column corresponding to bag $t$ has one entry for each
ancestor edge, namely $\operatorname{depth}_{T_i}(t)$ entries.
Consequently,
\[
 \norm{A_i}_2^2\le\norm{A_i}_F^2
       =\sum_{t\in T_i}\operatorname{depth}_{T_i}(t)
       \le\frac{m_i(m_i-1)}2\le G.
\]
The depth sum bound follows by ordering parents before children: the
vertex in position $j+1$ has depth at most $j$.  Applying this operator
bound to the coordinate error vector and summing over coordinates gives
\eqref{eq:inexact-gradient-incidence}.  The per-bag vector bound gives
$\sum_t\norm{e_t}^2\le N\gamma^2h_j^2$.  The componentwise condition
implies that vector bound because $|V_t|\le p$.
\end{proof}

There is no branching factor in either lemma.  Errors introduced while
forming subtree sums must be included in $\nu$ if those sums are
inexact.  Exact sums of rational gradient approximations avoid that
additional error.  No gradient accuracy is needed for slope
perturbations when $k=1$.

\subsection{Contraction and certified termination}

\begin{theorem}[Inexact regridding]
 \label{thm:inexact-main}
Use fresh partitions at $h_j=s_0 2^{-j}$ and center $x^{(j-1)}$, with
$\theta$ satisfying \eqref{eq:regridding-theta}.  Suppose fixed
constants $d,\zeta,\omega\ge0$ give, at every stage,
\begin{equation}
 \delta_{t,j}\le d h_j^2,\qquad
 \nu_j\le\zeta\sqrt N h_j,\qquad
 F(x^{(j)})\le U_j\le F(x^{(j)})+\omega Nh_j^2.
 \label{eq:inexact-budgets}
\end{equation}
Here $x^{(j)}$ is the consistent point of the backtracked configuration.
Maintain $\UB_j=\min_{0\le l\le j}U_l$, optionally also including a
certified feasible initial upper value.  Set
\begin{equation}
 K=C+K_s+d+\omega,\qquad
 B_{\mathrm{in}}=\max\left\{p,\frac{8K}{3g}\right\}.
 \label{eq:inexact-constants}
\end{equation}
Then, starting from any $x^{(-1)}\in X$, every stage satisfies
\begin{equation}
 \norm{x^{(j)}-x^*}^2\le B_{\mathrm{in}}Nh_j^2,
 \qquad
 0\le\UB_j-\LB_j\le gB_{\mathrm{in}}Nh_j^2.
 \label{eq:inexact-stage-bounds}
\end{equation}
For $\eps>0$, the gap test stops by stage
\[
 J=\max\left\{0,\ceil{\log_2\left(
               s_0\sqrt{gB_{\mathrm{in}}N/\eps}\right)}\right\}.
\]
The final number of leaves and cells is at most
$2N(4/\theta)^p(J+1)$.  Through stage $J$, the number of created boxes
is at most $N(4/\theta)^p(J+1)(J+2)$, and the number of local convex
oracle calls is at most
\[
 N3^q(4/\theta)^p\frac{(J+1)(J+2)(2J+3)}6.
\]
\end{theorem}

\begin{proof}
Every stage supplies a valid lower certificate by
\cref{prop:inexact-dp}, and every maintained upper value comes from a
feasible point.  Thus $\LB_j\le f^*\le\UB_j$.
Apply \eqref{eq:regridding-error} to the selected configuration with
the exact current-center slopes, then use
\eqref{eq:inexact-backtrack} and
\eqref{eq:inexact-slope-perturbation}.  Since $\eta=g/20$, their
combined squared-distance allowance is $\eta+\eta/4=g/16$, giving
\[
 F(x^{(j)})-\LB_j
    \le\frac g{16}\norm{x^{(j)}-x^{(j-1)}}^2
                      +(C+K_s+d)Nh_j^2.
\]
Put $e_j=\norm{x^{(j)}-x^*}^2$.  Quadratic growth, lower-bound
validity, and the inequality
$\norm{x^{(j)}-x^{(j-1)}}^2\le2(e_j+e_{j-1})$ imply
$ge_j\le g(e_j+e_{j-1})/8+KNh_j^2$.  Rearranging gives
\begin{equation}
 e_j\le\frac17e_{j-1}+\frac{8K}{7g}Nh_j^2.
 \label{eq:inexact-contraction}
\end{equation}
For $j\ge1$, an induction using $h_{j-1}=2h_j$ closes because
\[
 \frac47B_{\mathrm{in}}+\frac{8K}{7g}
       \le B_{\mathrm{in}}.
\]
For stage zero, $e_{-1}\le n s_0^2\le pNh_0^2
\le B_{\mathrm{in}}Nh_0^2$, and
\eqref{eq:inexact-contraction} yields the stronger
$e_0\le4B_{\mathrm{in}}Nh_0^2/7$.  This proves the distance bound.

For $j\ge1$, the two distance bounds imply
$\norm{x^{(j)}-x^{(j-1)}}^2\le10B_{\mathrm{in}}Nh_j^2$.
At stage zero the bound with coefficient $4$ suffices.  Since
$\UB_j\le U_j$, the evaluation allowance gives
\[
 \UB_j-\LB_j
   \le\left(\frac{5g}{8}B_{\mathrm{in}}+K\right)Nh_j^2
   \le gB_{\mathrm{in}}Nh_j^2,
\]
where the last step uses $K\le3gB_{\mathrm{in}}/8$.  The definition
of $J$ makes this upper bound at most $\eps$ at stage $J$.
At the stopping stage, $F(x_{\mathrm{inc}})-f^*\le\eps$ for the
feasible point associated with the maintained upper value.

The partitions and incidence lists are those of exact regridding.
The stage box bound from \cref{thm:regridding-main} and the incidence
bound of \cref{prop:regridding-incidence} therefore apply unchanged.
Summing $j+1$ and $(j+1)^2$ for $0\le j\le J$ gives the stated
created-box and local-call counts.  Local numerical iterations are not
included in these oracle-call counts.
\end{proof}

The allocation $\eta/4$ in the slope estimate preserves the strict
margin needed when widths halve: the coefficient $1/7$ in
\eqref{eq:inexact-contraction} is smaller than $1/4$.  Certified
lower and upper bounds make the stopping test sound even at grading
ratios that fail the convergence conditions.

\begin{corollary}[Unknown constants]
 \label{cor:inexact-search}
Fix positive local, gradient, and evaluation budgets $d,\gamma,\omega$,
and assume every requested certified oracle call terminates.  The
operation-budget search of \cref{prop:regridding-search} applies to
inexact regridding without knowledge of $g,M,A$: in round $r$, run each
$\theta=2^{-\mu}$, $1\le\mu\le r$, from the initial point, with a
budget of $2^r$ counted operations per trial and no separate stage cap.
Some trial terminates with a certified gap at most $\eps$.

More precisely, let $\mu_*$ be a sufficient grading index and let
$W_*\ge1$ bound the counted work of its completed run.  Put
$W=\max\{W_*,2^{\mu_*},2\}$ and $R=\ceil{\log_2 W}$.
The total counted work through termination is at most $2R2^R$, hence
$O(W\log W)$.  If box emission is counted, the returned certificate
contains at most $2^R$ boxes.
\end{corollary}

\begin{proof}
Request bag-gradient error $\gamma h_j$; \cref{lem:inexact-gradients}
supplies $\zeta=\sqrt G\,\gamma$ for the analysis.  None of these
requests depends on $g,M,A$.  A sufficiently small dyadic $\theta$
satisfies \eqref{eq:regridding-theta}, so the corresponding run stops
by \cref{thm:inexact-main}.  Its counted work is finite.  For example,
under the same unit-cost bag-oracle and arithmetic accounting as
\cref{prop:regridding-search}, it is
$O(pN3^q(4/\theta)^p(J+1)^3)$; oracle-internal work is excluded.
Enforce each trial budget before its next counted operation.  Round
$R$ includes $\mu_*$ and permits at least $W_*$ operations for that
trial, so that trial completes unless an earlier valid gap test has
already stopped the search.  The work bound follows from
$\sum_{r=1}^R r2^r\le2R2^R$, and $2^R<2W$.

The first successful trial may have a different grading ratio and
stage from the sufficient run.  Its guarantee is the certified gap
test.  Its actual stage has at most
$2N(4/\theta)^p(j+1)$ boxes, and its box emissions also fit within
its operation budget, which is at most $2^R$.  No final-box estimate
using the sufficient run's $\theta$ is asserted for this returned
trial.
\end{proof}

\subsection{Exact faces and rational reconstruction}

The feasibility requirement in \eqref{eq:inexact-oracle} also applies
when a leaf and a separator cell only touch.  A numerical displacement
of any size can leave a fixed face.  The following construction isolates
this issue from objective accuracy.

\begin{proposition}[Rational feasible points]
 \label{prop:inexact-reconstruction}
Assume the original box, initial center, widths, and grading ratio have
rational representations, and shell endpoints are formed by rational
arithmetic.  A local domain $C_{t,B,D}$ is a rational box, possibly
lower dimensional.  Its fixed coordinates can be removed exactly.
Rational local oracle points give a rational consistent point in $X$,
and therefore a rational center and rational partitions at the next
stage.

For quantitative reconstruction, suppose $q=q_{t,B,D}$ is
$L_q$-Lipschitz on its relative box, and a feasible point $v_0$ has a
certified upper value $u_0$ with
$q(v_0)\le u_0$ and $u_0-\ell\le\delta/2$.
Let $K_q\in\Q$ satisfy $K_q\ge L_q\sqrt p$.  If a rational feasible
point $v$ has $\abs{v_i-v_{0,i}}\le\tau$ in every coordinate, then
\begin{equation}
 q(v)\le u_0+K_q\tau=:u_{\mathrm{new}}.
 \label{eq:inexact-rounding}
\end{equation}
For $\delta>0$ and $K_q>0$, rational $\tau>0$ with
$K_q\tau\le\delta/2$ ensures
$u_{\mathrm{new}}-\ell\le\delta$.  If $K_q=0$, reconstruction has
no objective penalty.
\end{proposition}

\begin{proof}
For each coordinate the intersection interval is
$[\max\{L_i^B,L_i^D\},\min\{U_i^B,U_i^D\}]$, with the own-cell
endpoints omitted outside the separator.  Exact endpoint comparisons
determine emptiness.  Equal endpoints fix that coordinate to the common
rational value; eliminate it before solving for the remaining
coordinates.  Restricting the objective preserves continuity and
convexity.  If no coordinates remain, the task is certified evaluation
at the single rational point.  The consistent point selects one of the
feasible bag coordinates for each original variable, so it is rational
and lies in its original interval.  Rational shell arithmetic then
preserves the claim at the next stage.

To construct $v$, obtain a rational approximate vector $w$ with
certified coordinate errors at most $\tau$ from $v_0$ and clip each
coordinate exactly to its intersection interval.  Fixed coordinates
are set to their exact endpoint.  Since $v_0$ lies in the interval,
clipping cannot increase its coordinate distance from $v_0$.
Thus $v$ is rational and feasible, with
$\norm{v-v_0}\le\sqrt p\,\tau$.  The Lipschitz estimate gives
$q(v)\le q(v_0)+L_q\sqrt p\,\tau\le u_0+K_q\tau$.
The remaining assertions follow by adding this allowance to the
initial certified gap.
\end{proof}

Exact clipping certifies feasibility, not the objective gap.  Without
a supplied quantitative modulus, the oracle must certify the final
point and its gap directly.  A fresh evaluation after reconstruction
must fit within the remaining gap budget.  Continuity ensures that
rational feasible points approximate a minimum on a rational relative
box at every positive tolerance, since such points are dense there.
It does not supply a procedure or cost bound for finding and certifying
them.  Rational gradient approximations and exact rational subtree sums
also give rational slopes; rational lower bounds give rational message
intercepts.

Zero budgets require exact oracles and an admissible exact
representation.  They need not admit rational points, even for rational
convex polynomial data.  On $[1,2]$, the function
$q(x)=(x^2-2)^2$ is strictly convex because $q''(x)=12x^2-8\ge4$.
Its unique minimizer is $\sqrt2$.  With $\delta=0$,
\eqref{eq:inexact-oracle} forces the returned point to equal that
minimizer, so no rational point can satisfy the contract.  This is why
oracle realizability and termination remain assumptions.  Irrational
fixed domain coordinates likewise require an appropriate exact
representation.

For fixed positive budgets, the requested absolute objective accuracy
is proportional to $h_j^2$ and the gradient accuracy to $h_j$, so their
precision requests grow linearly with $j$, apart from scale terms.
This observation does not bound the bit lengths of all certificate
entries or the internal work of the oracles.  The certificate size
counts leaves and cells; a serialized proof may include separate local
lower-bound witnesses for every touching pair and may consequently
contain more proof objects.  Bit-complexity conclusions require an
explicit representation and certified routines with their own precision
and runtime bounds.
